\documentclass[a4paper,12pt]{article}
\usepackage[utf8]{inputenc}
\usepackage[russian]{babel}
\usepackage{amsmath, amssymb}

\begin{document}

\section{Пример анализа циклического кода}

Рассмотрим двоичный циклический код длины $n = 7$, заданный порождающим многочленом
\[
g(x) = 1 + x^2 + x^3.
\]

Степень многочлена равна $3$, следовательно, размерность кода определяется как
\[
k = n - \deg g(x) = 7 - 3 = 4.
\]

Таким образом, получаем код с параметрами $(7,4)$.

Все кодовые слова формируются по правилу
\[
c(x) = m(x)g(x) \mod (x^7 - 1),
\]
где информационный многочлен $m(x)$ имеет степень строго меньше $4$.

Рассмотрим несколько примеров:

\begin{itemize}
    \item $m(x) = 0 \Rightarrow c(x) = 0000000$;
    
    \item $m(x) = 1 \Rightarrow c(x) = g(x) = 1011000$;
    
    \item $m(x) = x \Rightarrow c(x) = xg(x) = 0110100$;
    
    \item $m(x) = x^2 \Rightarrow c(x) = 0011010$;
    
    \item $m(x) = x^3 \Rightarrow c(x) = 0001101$.
\end{itemize}

Остальные кодовые слова получаются аналогично за счёт линейных комбинаций.

Минимальное расстояние данного кода равно $d_{\min} = 3$, что следует из анализа весов ненулевых кодовых слов. Следовательно, рассматриваемый код совпадает с классическим кодом Хэмминга $(7,4)$.

\section{Двойственный код}

Пусть задан циклический код $C$ длины $n$ над полем $\mathbb{F}_q$ с порождающим многочленом $g(x)$. Тогда можно определить проверочный многочлен:
\[
h(x) = \frac{x^n - 1}{g(x)}.
\]

Двойственный код $C^\perp$ также является циклическим. Его порождающий многочлен получается из $h(x)$ изменением порядка коэффициентов:
\[
\tilde{g}(x) = x^{\deg h} h(x^{-1}).
\]

Иначе говоря, если
\[
h(x) = h_0 + h_1 x + \dots + h_{n-k} x^{n-k},
\]
то
\[
\tilde{g}(x) = h_{n-k} + h_{n-k-1}x + \dots + h_0 x^{n-k}.
\]

Таким образом, операция двойственности сохраняет цикличность кода.

\section{Общий способ построения циклических кодов}

Построение циклического кода можно описать следующим алгоритмом:

\begin{enumerate}
    \item Разложить многочлен $x^n - 1$ на неприводимые множители над полем $\mathbb{F}_q$.
    
    \item Выбрать делитель $g(x)$ степени $n - k$ — он и будет порождающим многочленом.
    
    \item Сформировать кодовые слова как
    \[
    c(x) = m(x)g(x) \mod (x^n - 1),
    \]
    где $\deg m(x) < k$.
    
    \item При необходимости построить систематическую форму, используя деление:
    \[
    c(x) = m(x)x^{n-k} - \big(m(x)x^{n-k} \bmod g(x)\big).
    \]
\end{enumerate}

Данный подход гарантирует линейность и циклическую структуру кода.

\section{Основные алгебраические структуры}

\subsection{Группы}

Группа — это множество $G$ с операцией, обладающей следующими свойствами:
замкнутость, ассоциативность, наличие нейтрального элемента и обратимых элементов.

Если дополнительно выполняется коммутативность, группа называется абелевой.

\subsection{Смежные классы и факторизация}

Для подгруппы $H \subset G$ можно определить смежные классы:
\[
gH = \{gh \mid h \in H\}.
\]

Они разбивают группу на непересекающиеся множества одинаковой мощности. Число таких классов называется индексом подгруппы.

\subsection{Теорема Лагранжа}

Если группа конечна, то порядок любой её подгруппы делит порядок всей группы:
\[
|G| = |H| \cdot [G : H].
\]

\subsection{Фактор-группа}

Множество смежных классов образует новую группу $G/H$, называемую фактор-группой.

\section{Кольца и поля}

Кольцо — это множество с операциями сложения и умножения, где сложение образует абелеву группу, а умножение ассоциативно и дистрибутивно относительно сложения.

Поле — это коммутативное кольцо с единицей, в котором каждый ненулевой элемент обратим.

Примеры: $\mathbb{Q}$, $\mathbb{R}$, $\mathbb{F}_p$.

\section{Многочлены и расширения полей}

Кольцо многочленов $\mathbb{F}_q[x]$ состоит из всех полиномов с коэффициентами из $\mathbb{F}_q$.

Если взять факторкольцо
\[
\mathbb{F}_q[x]/(p(x)),
\]
где $p(x)$ — неприводимый многочлен степени $m$, то получается поле из $q^m$ элементов.

\subsection{Пример}

Пусть $p(x) = 1 + x^2$ над $\mathbb{F}_2$. Тогда элементы факторкольца имеют вид:
\[
a_0 + a_1 x, \quad a_i \in \mathbb{F}_2,
\]
и всего таких элементов $4$.

Из соотношения $x^2 \equiv 1$ следует правило умножения. Например:
\[
x \cdot x = 1,
\]
поэтому элемент $x$ является обратным самому себе.

\end{document}